% Independent preview; original geometry and numeric relationships retained.
% 教学构造：两个面板均取 g(x)=x_1-2，规范间隔边界为 x_1=1 与 x_1=3。
\begin{tikzpicture}[foundation picture,x=1cm,y=1cm,font=\figurefont\footnotesize]
  \foreach \shift/\title in {0/硬间隔,5.6/软间隔允许违反约束} {
    \begin{scope}[xshift=\shift cm]
      \fill[FigureYellowFill] (1,0) rectangle (3,4);
      \draw[foundation axis] (0,0) -- (4.35,0) node[right] {$x_1$};
      \draw[foundation axis] (0,0) -- (0,4.2) node[above] {$x_2$};
      \draw[FigureStroke,line width=1.1pt] (2,0) -- (2,4);
      \draw[FigureBlue,dashed,line width=1pt] (1,0) -- (1,4);
      \draw[FigureBlue,dashed,line width=1pt] (3,0) -- (3,4);
      \node[foundation heading] at (2,4.8) {\title};
      \node[below,font=\footnotesize] at (1,-0.1) {$g=-1$};
      \node[below,font=\footnotesize] at (2,-0.1) {$g=0$};
      \node[below,font=\footnotesize] at (3,-0.1) {$g=1$};
      \foreach \x/\y in {0.5/2.2,0.7/3.3,1/1.1}
        \fill[FigureBlue] (\x,\y) circle (2.4pt);
      \foreach \x/\y in {3/2.6,3.6/1.9,3.7/3.4}
        \fill[LancetGreen] (\x-0.08,\y-0.08) rectangle ++(0.16,0.16);
    \end{scope}
  }
  \draw[FigureStroke] (1,1.1) circle (4.5pt);
  \draw[FigureStroke] (3,2.6) circle (4.5pt);
  \draw[{Stealth[length=1.5mm]}-{Stealth[length=1.5mm]},FigureStroke]
    (1,3.9) -- node[above,font=\footnotesize] {$2/\lVert\bm{w}\rVert_2$} (3,3.9);
  \begin{scope}[xshift=5.6cm]
    \fill[LancetGreen] (2.5-0.08,0.9-0.08) rectangle ++(0.16,0.16);
    \node[above left] at (2.5,0.9) {A};
    \fill[LancetGreen] (1.5-0.08,2.9-0.08) rectangle ++(0.16,0.16);
    \node[above left] at (1.5,2.9) {B};
    \draw[FigureRed,line width=1.5pt] (2.5,0.9) -- (3,0.9)
      node[midway,below,text=FigureInk,inner sep=1pt] {$\xi_A=0.5$};
    \draw[FigureRed,line width=1.5pt] (1.5,2.9) -- (3,2.9)
      node[midway,above,text=FigureInk,inner sep=1pt] {$\xi_B=1.5$};
  \end{scope}
  \fill[FigureBlue] (2.7,-1) circle (2.4pt);
  \node[right] at (2.85,-1) {负类};
  \fill[LancetGreen] (5.9,-1.08) rectangle ++(0.16,0.16);
  \node[right] at (6.15,-1) {正类};
\end{tikzpicture}
